\section{Method}
\label{sec:method}

Figure~\ref{fig:pipeline} summarises the pipeline. All components run locally on a 4\,GB laptop GPU.

\subsection{Evidence units and adaptive retrieval}
Sources (PDF, DOCX, text, web pages) are split into sentences; each sentence is an \emph{evidence unit} with a stable identifier, page, section and source, so every statement can be cited.
Each unit $u$ is scored for a query $q$ by
\begin{equation}
  s(u,q) = \lambda \cos(\mathbf{e}_u,\mathbf{e}_q) + (1-\lambda)\,\widetilde{\mathrm{BM25}}(u,q) + \beta\,\mathrm{cov}(u,q)\,\mathbb{1}[\tau(q)\in\phi(u)],
\end{equation}
with MiniLM sentence embeddings \citep{reimers2019sbert,wang2020minilm}, BM25 \citep{robertson2009bm25} normalised by the score of an ideal match, and an intent bonus that is non-zero only if the unit has the evidence type the question asks for (causal, comparative, numeric, limitation; $\tau(q)\in\phi(u)$), scaled by the fraction $\mathrm{cov}(u,q)$ of the question's key terms the unit contains.
The retriever keeps up to $k_0$ units scoring at least $\rho\,s_{\max}$ (trimming), then tests \emph{sufficiency}: best score, key-term coverage, presence of every named entity of the question, and the presence of the requested evidence type.
If insufficient, it expands in up to two rounds (missing-term search, neighbouring sentences, relaxed cutoff).
We use $\lambda{=}0.65$, $\beta{=}0.08$, $k_0{=}4$, $\rho{=}0.7$ throughout.

\subsection{Answerability gate}
Retrieved evidence can be topically relevant without answering the question.
An extractive reader trained on SQuAD~2.0 reads the evidence and returns its best span together with the margin between the best span score and the null (``no answer'') score.
The system abstains when the margin is negative (the reader's own decision; not tuned) and otherwise leads the answer with the sentence containing the span.
Answers are extractive: every answer sentence is a cited source sentence.

\subsection{Claim-level verification}
Answers (or any text, e.g.\ an LLM draft) are split into claims conservatively: sentences, split further only at clause boundaries with self-contained halves; attribution wrappers (``According to the passages, \dots'') and meta statements are removed.
For a claim $c$, the verifier gathers candidate premises: the $m$ most relevant single sentences, adjacent-sentence windows, and one \emph{multi-sentence} premise built from the top-$k$ sentences in document order (for claims that aggregate facts).
A cross-encoder NLI model \citep{he2023debertav3} returns $P(\text{entail}), P(\text{contra}), P(\text{neutral})$ for each premise, and the claim is labelled
\textsc{Supported} if a relevant premise entails it (and all its numbers occur in the premise),
\textsc{Contradicted} if a closely relevant \emph{single} sentence (or a window more relevant than every single sentence) contradicts it,
\textsc{Uncertain} if support and contradiction are of comparable strength (relevance\,$\times$\,probability within a margin) or the NLI probabilities are inconclusive,
\textsc{Partially supported} if only a component of a causal claim (``A because B'') is supported,
and \textsc{Insufficient evidence} otherwise.
Missing evidence is never labelled contradicted.
Multi-sentence premises are support-only because NLI models produce spurious contradictions on long multi-entity premises.

\subsection{Budget-aware verification}
Checking every candidate premise costs $\approx$10--15 NLI calls per claim.
The claims of an answer instead share a budget $B = \max(B_{\min}, b\cdot n)$ of NLI calls.
Premises of each claim are ordered by relevance.
Every claim first receives one step of two checks; afterwards the scheduler picks the unresolved claim maximising
\begin{equation}
  \frac{\pi_c\,(1-d_c)}{(1+\alpha\,a_c)(1+\gamma\,\ell_c)},\qquad
  \pi_c = 0.7\,(1-r_c) + 0.3\,h_c,
\end{equation}
where $r_c$ is the relevance of the claim's best candidate (so $1-r_c$ is its evidence gap), $h_c$ a hedging-based uncertainty, $d_c$ the current decisiveness (best relevant entailment or contradiction probability), $a_c$ the number of steps taken and $\ell_c$ the streak of steps whose evidence gain (increase of $d_c$) was below $0.05$.
A claim stops consuming budget as soon as it is \textsc{Supported} or \textsc{Contradicted} (early stopping), and part of the budget is reserved for the components of causal claims.
Claims the budget never reached are reported as \textsc{Uncertain}, never as supported.
We compare this priority scheduler with round-robin scheduling (same early stopping) and a fixed per-claim split (no early stopping).

\subsection{Revision}
Given claim verdicts, the answer is rewritten: supported claims are kept with a citation, contradicted claims are replaced by a correction quoting the source, only the established part of partially supported claims is kept, uncertain claims are marked, and unsupported claims are removed.

\subsection{Fine-tuning the verifier on RAGTruth}
To adapt the verifier to LLM outputs we fine-tune the NLI cross-encoder on RAGTruth \emph{train}.
Every claim extracted from a training response becomes an example whose premise is the top-5 source sentences retrieved by our retriever.
Labels keep NLI semantics so that the fine-tuned model is a drop-in replacement: faithful $\rightarrow$ entailment, \emph{conflict} $\rightarrow$ contradiction, \emph{baseless} $\rightarrow$ neutral.
We hold out 10\% of the training \emph{sources} for validation, subsample faithful claims (35\%) to reduce imbalance, freeze the word embeddings, and train for two epochs with fp16 and gradient checkpointing (peak 2.1\,GB of GPU memory).
The RAGTruth test split is used only for the final evaluation.
